\documentclass[12pt, a4paper]{article}

% ── Packages ──────────────────────────────────────────────
\usepackage[utf8]{inputenc}
\usepackage[T1]{fontenc}
\usepackage{mathptmx}                    % Times New Roman
\usepackage[margin=1in]{geometry}
\usepackage{graphicx}
\usepackage{booktabs}                     % Professional tables
\usepackage{multirow}                     % Merged table cells
\usepackage{caption}
\usepackage{subcaption}                   % Side-by-side figures
\usepackage{amsmath}
\usepackage{hyperref}
\usepackage{float}
\usepackage{enumitem}
\usepackage{setspace}
\usepackage{natbib}
\usepackage{xcolor}

\hypersetup{
    colorlinks=true,
    linkcolor=blue!60!black,
    citecolor=blue!60!black,
    urlcolor=blue!60!black
}

\onehalfspacing

% ── Title ─────────────────────────────────────────────────
\title{
    \vspace{-1cm}
    \Large\textbf{SkillCraft1: Multinomial Skill Classification\\Using Machine Learning}
}
\author{
    William Chen\\
    \textit{University of California, Irvine}\\
    Instructor: Prof.\ Sam Behseta\\
    \textit{DATA 200BP: Intermediate Probability and Statistical Theory II}
}
\date{}

% ══════════════════════════════════════════════════════════
\setlength{\parindent}{0pt}
\setlength{\parskip}{8pt}
\begin{document}
\maketitle
\thispagestyle{empty}

% ── Abstract ──────────────────────────────────────────────
\begin{abstract}
This project extends the work of Thompson et al.\ (2013), who used pairwise binary
classification with Conditional Inference Forest to study variable importance across
StarCraft~II skill levels. While their approach provided insights into how predictors
of expertise change across leagues, it did not address the full complexity of multi-class
skill prediction or the class imbalance inherent in the SkillCraft1 dataset. This study
reformulates the problem as a 6-class multinomial classification task, predicting player
league rankings from Bronze~(1) to Masters~(6).

We evaluate three machine learning models: Multinomial Logistic Regression, Decision Tree, and Random Forest. Each model is trained under three class imbalance strategies (Baseline, Class Weights, and SMOTE combined with Tomek Links), and performance is assessed using overall accuracy and minority class recall on a held-out test set.

Results show that Random Forest with Class Weights achieves the best overall balance,
with an accuracy of 0.4070 and a meaningful improvement in minority class recall
compared to the baseline. In a direct pairwise comparison replicating Thompson et al.'s
setup, our best model outperforms their Conditional Inference Forest results in 3 out
of 4 league pairs. These findings suggest that class weighting is an effective and
practical strategy for handling imbalance in multi-class skill classification tasks.
\end{abstract}

\newpage
\tableofcontents
\newpage

% ══════════════════════════════════════════════════════════
\section{Introduction}
% ══════════════════════════════════════════════════════════

\subsection{Nature of the Project}

Cognitive science has long studied expertise by comparing experts with novices. Thompson
et al.\ (2013) challenged the assumptions underlying such comparisons by analyzing
StarCraft~II telemetry data from 3,395 players across competitive leagues. They discovered that variable importance is not static: APM dominates at lower skill levels, but at higher levels, attentional efficiency metrics become more predictive.

Their method used pairwise binary classifiers built with Conditional Inference Forests,
each comparing players two leagues apart. This design was deliberate: neighboring leagues
are not cleanly separable because league placement does not perfectly reflect objective
skill. They found that classifiers performed poorly on adjacent leagues and had
significantly more success when the distance between classes was at least two.

StarCraft~II is well-suited for this type of study. Unlike chess, an average RTS game
contains \textasciitilde1,635 moves, providing dense behavioral data in a naturalistic
setting. The SkillCraft1 dataset contains 20 behavioral and demographic variables, making
it the largest expertise study of its kind.

\subsection{Aims}

This project extends Thompson et al.'s pairwise approach in two directions. First, we
replicate their pairwise binary classifiers using our best-performing configuration to
benchmark our methodology against their Conditional Inference Forest results. Second, we develop a direct 6-class multinomial prediction model that simultaneously classifies players across all skill levels, a more challenging problem that requires handling severe class imbalance and non-linear decision boundaries.

To address class imbalance, three strategies are compared: baseline classification,
class weight adjustment, and SMOTE combined with Tomek Links. This allows us to
investigate whether imbalance-aware training improves prediction of minority classes.

\subsection{Goals}

The primary goal of this project is to develop an effective 6-class multinomial
classifier for Leagues 1 through 6 that balances overall accuracy with minority class
recall. To achieve this, we compare three machine learning models: Multinomial Logistic Regression, Decision Tree, and Random Forest. Each model is trained under three imbalance strategies (Baseline, Class Weights, and SMOTE combined with Tomek Links). Additionally, we replicate Thompson et al.'s pairwise binary classification setup to benchmark our best-performing model against their Conditional Inference Forest results.

A secondary goal is to examine whether demographic features such as Age and HoursPerWeek
contribute meaningfully to skill prediction beyond in-game behavioral metrics.


% ══════════════════════════════════════════════════════════
\section{Literature Review}
% ══════════════════════════════════════════════════════════

\subsection{Skill Classification in Competitive Gaming}

Skill assessment in competitive games bridges expertise research with machine learning.
Unlike laboratory settings, games provide naturalistic data where players compete under
identical conditions. Real-time strategy games like StarCraft~II are particularly rich:
they demand simultaneous management of multiple objectives and high-frequency
decision-making, making them ideal for studying expertise across a continuous skill
spectrum.

\subsection{Thompson et al.\ (2013): SkillCraft1 Baseline}

Thompson et al.\ (2013) introduced the SkillCraft1 dataset, comprising 3,395 players
across 7 leagues with 20 behavioral and demographic variables. Their primary objective
was to examine how variable importance changes across skill levels, rather than to
maximize classification accuracy. Using pairwise binary classifiers built with
Conditional Inference Forests (ntree = 1000, mtry = 5), they compared leagues two
levels apart to avoid the separability issues inherent in adjacent league comparisons.

Their analysis revealed that adjacent leagues are difficult to separate because league
placement does not perfectly reflect objective skill. They also found that variable
importance shifts across the skill spectrum: APM predicts skill at lower levels, while
attentional efficiency metrics such as Action Latency and Gap Between PACs become
dominant at higher levels. Grandmaster players were excluded due to their small sample
size (35 players), and 55 professional players were added from publicly available
tournament replays to form the highest skill group.

This project extends their work by replicating the pairwise comparison using RF + Class Weights (a standard Random Forest with class weighting, as opposed to their Conditional Inference Forest) and by attempting direct 6-class multinomial prediction, a harder problem that Thompson et al.\ did not pursue due to the separability limitations they observed.

\subsection{Class Imbalance Handling Methods}

Class imbalance occurs when certain classes are significantly smaller than others,
causing models to favor majority classes and underperform on minority ones. Three
common strategies for addressing this issue are evaluated in this study.

SMOTE generates synthetic minority samples by interpolating between existing
observations and their nearest neighbors. This increases representation without
duplication, but may produce samples that do not reflect realistic data patterns.

Tomek Links remove majority class samples that are nearest neighbors of minority
class samples, cleaning the decision boundary. This is often combined with SMOTE
to provide both augmentation and boundary clarity.

Class Weighting adjusts the loss function to penalize misclassification of minority
samples more heavily, typically using inverse frequency weighting. Unlike resampling
methods, it preserves the original data distribution entirely.

\subsection{Machine Learning Models for Classification}

Multinomial Logistic Regression produces linear decision boundaries and interpretable
coefficients, but may underperform when feature-class relationships are nonlinear.

Decision Trees partition the feature space using binary splits, naturally handling
nonlinearity. However, they are prone to overfitting and require pruning to generalize.

Random Forests aggregate many decision trees through bagging and random feature
selection, reducing overfitting and improving robustness at the cost of interpretability.

These three models are compared across three imbalance strategies to identify the
most effective combination for 6-class skill classification.


% ══════════════════════════════════════════════════════════
\section{Data Introduction}
% ══════════════════════════════════════════════════════════

\subsection{SkillCraft1 Dataset Overview}

The SkillCraft1 dataset contains 3,395 StarCraft~II player observations across 8
leagues. After removing League~8 (systematic missing data in all demographic features)
and League~7 (Grandmaster, excluded per Thompson et al.\ due to small sample size),
the analysis uses 3,340 observations across 6 leagues: Bronze, Silver, Gold, Platinum,
Diamond, and Masters.

The dataset is severely class-imbalanced, with Bronze players comprising only about
5\% of observations while Platinum and Diamond together account for nearly half.

\subsection{Data Features \& Characteristics}

The dataset contains 20 variables: 3 demographic (Age, HoursPerWeek, TotalHours) and
17 behavioral, capturing in-game actions such as APM, hotkey usage, map control,
decision-making efficiency, and unit production, all normalized per minute of gameplay.


% ══════════════════════════════════════════════════════════
\section{Exploratory Data Analysis}
% ══════════════════════════════════════════════════════════

\subsection{Summary Statistics}

From the summary, \texttt{Age}, \texttt{HoursPerWeek}, and \texttt{TotalHours} contain
missing values (NA). Additionally, \texttt{TotalHours} shows an abnormally large maximum
value of 1,000,000, suggesting the presence of an outlier that will need to be addressed
in preprocessing.

\subsection{Class Distribution}

The distribution of players across leagues is highly imbalanced (Figure~\ref{fig:class-dist}).
Platinum (League~4) and Diamond (League~5) account for the majority of observations,
while Grandmaster (League~7) and Professional (League~8) have very few samples.
This imbalance will be addressed in Section~5 and Section~6.

\begin{figure}[H]
    \centering
    \includegraphics[width=0.7\textwidth]{figures/class_distribution.png}
    \caption{Player distribution by league.}
    \label{fig:class-dist}
\end{figure}

\subsection{Feature Distribution by League}

Four features are examined to assess their relationship with league rank
(Figure~\ref{fig:feature-boxplots}). Age shows a stable median around 20--22 across all
leagues, suggesting age alone is not a strong indicator of skill level. MinimapAttacks
shows narrow boxes close to zero across most leagues, though higher leagues show slightly
wider distributions. TotalHours contains an extreme outlier (1,000,000 hours) that
dominates the y-axis, confirming the need for outlier capping in preprocessing.
WorkersMade shows the highest median at League~6 rather than League~7 or~8, suggesting
worker production does not linearly increase with rank.

\begin{figure}[H]
    \centering
    \includegraphics[width=0.85\textwidth]{figures/feature_boxplots.png}
    \caption{Feature distributions across leagues for Age, MinimapAttacks, TotalHours,
    and WorkersMade.}
    \label{fig:feature-boxplots}
\end{figure}

\subsection{Data Quality Issues}

TotalHours contains an extreme outlier of 1,000,000 hours. After capping at the 99th
percentile (2,520 hours), the distribution shows a right-skewed pattern with most
players under 1,000 hours. Actual outlier handling is applied in Section~5.

All 55 League~8 observations are missing Age, HoursPerWeek, and TotalHours (100\%),
indicating systematic missingness, likely because professional players' data was
collected from tournament records rather than standard ladder profiles. In contrast,
League~5 has only 2--3 missing entries ($<0.3$\%), consistent with random gaps.
This supports excluding League~8 entirely rather than attempting imputation.

\begin{figure}[H]
    \centering
    \includegraphics[width=0.7\textwidth]{figures/totalhours_dist.png}
    \caption{Distribution of TotalHours after capping at the 99th percentile (2,520 hours).
    The distribution is right-skewed, with most players logging under 1,000 hours.}
    \label{fig:totalhours-dist}
\end{figure}

\subsection{Feature Correlations}

The correlation heatmap reveals several notable relationships with LeagueIndex
(Figure~\ref{fig:corr-heatmap}). APM (0.66) and AssignToHotkeys (0.53) show strong
positive correlations, indicating higher-ranked players have faster actions and more
hotkey usage. ActionLatency ($-0.67$) shows the strongest negative correlation,
confirming that lower reaction time is associated with higher league rank. APM and
ActionLatency are strongly negatively correlated ($-0.72$), suggesting potential
multicollinearity to be addressed in modeling.

\begin{figure}[H]
    \centering
    \includegraphics[width=0.6\textwidth]{figures/correlation_heatmap.png}
    \caption{Correlation heatmap of selected features with LeagueIndex. ActionLatency
    ($-0.67$) and APM (0.66) show the strongest associations with league rank.}
    \label{fig:corr-heatmap}
\end{figure}


% ══════════════════════════════════════════════════════════
\section{Data Preprocessing}
% ══════════════════════════════════════════════════════════

League~8 (Professional) is excluded due to 100\% systematic missing values in all
demographic features. League~7 (Grandmaster) is excluded per Thompson et al.\ (2013)
due to its small sample size and analytic difficulties. TotalHours is capped at the
99th percentile to remove extreme outliers while preserving the overall distribution.
The remaining missing values (5 observations in League~5) are imputed using
\texttt{rfImpute}, which uses Random Forest proximity to estimate missing values
iteratively over 6 iterations. The OOB error stabilizes around 58--59\%, indicating
convergence. The final dataset contains 3,340 observations across 6 leagues
(Bronze through Masters), split into 80\% training ($n=2{,}644$) and 20\% test
($n=661$) sets with a fixed random seed.


% ══════════════════════════════════════════════════════════
\section{Methods}
% ══════════════════════════════════════════════════════════

\subsection{Feature Selection and Tuning}

\paragraph{Multicollinearity Check (VIF).}
VIF analysis identifies four features with high multicollinearity: APM (41.05),
SelectByHotkeys (13.76), NumberOfPACs (14.96), and ActionsInPAC (9.42). These
overlap substantially with other action-based and PAC-related features. ActionLatency
(5.24) is also moderately inflated. All remaining features have VIF below 3.

\paragraph{Backward Elimination.}
Backward elimination iteratively removes features that contribute least to model fit based
on AIC. The final model retains 14 features, removing APM, MinimapRightClicks,
ComplexAbilitiesUsed, and Age.

\paragraph{Decay Tuning.}
The decay (L2 regularization) parameter is tuned using 5-fold cross-validation. Results
show that decay~$= 0$ achieves the highest accuracy (0.4217), meaning no regularization
is optimal. The final logistic regression model is built with decay~$= 0$.

\subsection{Class Imbalance Strategies}

As shown in Section~4.2, the dataset is heavily imbalanced. League~1 (Bronze) comprises
only 5.05\% of observations, corresponding to approximately 134 training samples after the
80/20 split. Three strategies are applied to address this imbalance, as described below.

\begin{enumerate}[itemsep=2pt]
    \item \textbf{Baseline}: No modification to the training data.
    \item \textbf{Class Weights}: Inverse-frequency weighting so smaller classes receive
          proportionally higher loss penalties, without modifying the data.
    \item \textbf{SMOTE + Tomek Links}: SMOTE oversamples minority classes via
          nearest-neighbor interpolation ($K=5$), then Tomek Links removes majority-class
          borderline samples. Applied only to the training set to prevent data leakage.
\end{enumerate}

\subsection{Machine Learning Models}

\paragraph{Multinomial Logistic Regression.}
Uses the 14 features selected via backward elimination (Age, APM, MinimapRightClicks,
and ComplexAbilitiesUsed removed). Trained with maximum likelihood estimation
(max iterations = 500). Decay parameter is tuned via 5-fold cross-validation,
with decay~$= 0$ selected as optimal.

\paragraph{Decision Tree.}
CART algorithm with all available features. Pruned based on cross-validation error
(\texttt{xerror}) to prevent overfitting.

\paragraph{Random Forest.}
Ensemble of 1,000 trees via bootstrap aggregating (bagging) with random feature
selection at each split. A proximity matrix is computed to enable distance-based
diagnostics. All features included, consistent with Thompson et al.\ (2013).

\subsection{Evaluation Metrics}

A unified evaluation function is applied across all three models. It generates a
confusion matrix on the test set and reports two metrics: overall accuracy and
recall for the minority class (League~1, Bronze).

% ══════════════════════════════════════════════════════════
\section{Results}
% ══════════════════════════════════════════════════════════

\subsection{Overall Model Comparison}

Table~\ref{tab:all-results} summarizes the performance of all nine model--strategy
combinations. RF + Class Weights achieves the best balance, with the highest accuracy
among imbalance-aware strategies (0.4070). SMOTE + Tomek consistently produces the
highest minority recall across all three models, but at a significant cost to overall
accuracy. Decision Tree is the weakest performer overall, with accuracy dropping below
0.30 under both Class Weights and SMOTE + Tomek. Logistic Regression Baseline achieves
the highest raw accuracy (0.4145), but its minority recall of 0.2727 indicates it largely
ignores the minority class. Given the class imbalance objective of this project, RF +
Class Weights is selected as the best overall model.

\begin{table}[H]
\centering
\caption{Performance comparison of all nine model--strategy combinations. Bold values
indicate the best result in each column.}
\label{tab:all-results}
\begin{tabular}{llcc}
\toprule
\textbf{Algorithm} & \textbf{Strategy} & \textbf{Accuracy} & \textbf{Minority Recall} \\
\midrule
\multirow{3}{*}{Logistic Regression}
    & Baseline        & \textbf{0.4145} & 0.2727 \\
    & Class Weights   & 0.3812          & 0.5152 \\
    & SMOTE + Tomek   & 0.3661          & 0.9091 \\
\midrule
\multirow{3}{*}{Decision Tree}
    & Baseline        & 0.3767          & 0.4545 \\
    & Class Weights   & 0.2874          & 0.3939 \\
    & SMOTE + Tomek   & 0.2844          & \textbf{0.9091} \\
\midrule
\multirow{3}{*}{Random Forest}
    & Baseline        & 0.4024          & 0.3333 \\
    & Class Weights   & 0.4070          & 0.3333 \\
    & SMOTE + Tomek   & 0.3933          & 0.6667 \\
\bottomrule
\end{tabular}
\end{table}

\subsection{Logistic Regression Results}

The baseline achieves an accuracy of 0.4145 but a minority recall of only 0.2727,
indicating most Bronze players are misclassified as Silver or Gold. Class Weights
improves minority recall to 0.5152 with only a marginal accuracy drop to 0.3812.
SMOTE + Tomek achieves the highest minority recall of 0.9091 but the lowest accuracy
(0.3661), as synthetic oversampling disrupts decision boundaries for mid-tier leagues.
Class Weights achieves the best tradeoff for this model.

\subsection{Decision Tree Results}

The baseline achieves an accuracy of 0.3767 with a minority recall of 0.4545.
ActionLatency is the primary splitting feature, with the entire tree built around three
thresholds of this variable. Notably, League~2 (Silver) appears as unused; the pruned
tree predicts zero samples as Silver.

\begin{figure}[H]
    \centering
    \includegraphics[width=0.85\textwidth]{figures/decision_tree.png}
    \caption{Pruned Decision Tree (baseline). The tree splits exclusively on ActionLatency
    thresholds; League~2 (Silver) is never predicted.}
    \label{fig:decision-tree}
\end{figure}

Unlike Logistic Regression, Class Weights does not improve minority recall for Decision
Trees; recall actually drops from 0.4545 to 0.3939, while accuracy falls substantially
to 0.2874. SMOTE + Tomek achieves the highest minority recall (0.9091) but at the largest
accuracy cost (0.2844). Decision Trees appear more sensitive to imbalance strategy than
Logistic Regression, with both strategies causing mid-tier league predictions to collapse
entirely.

\subsection{Random Forest Results}

Random Forest is the most stable model across all three imbalance strategies. The
baseline achieves an accuracy of 0.4024 with a minority recall of 0.3333. Class Weights
achieves a marginal accuracy improvement to 0.4070 with no change in minority recall,
suggesting the ensemble is already relatively robust to class imbalance through its
bagging mechanism. SMOTE + Tomek improves recall to 0.6667 with only a modest accuracy
drop to 0.3933. Unlike Decision Tree, no class collapses to zero under any strategy.

\paragraph{Variable Importance.}
The variable importance plot (Figure~\ref{fig:varimp}) confirms that ActionLatency and APM
are the two most influential predictors, consistent with Thompson et al.'s finding that
action speed dominates skill prediction. NumberOfPACs and TotalHours follow as the next
most important features. Demographic features such as Age and HoursPerWeek rank in the
lower half, suggesting that in-game behavioral metrics are more predictive of skill than
time investment alone. ComplexUnitsMade is the least important feature overall.

\begin{figure}[H]
    \centering
    \includegraphics[width=0.7\textwidth]{figures/variable_importance.png}
    \caption{Variable importance from Random Forest (SMOTE + Tomek). ActionLatency and APM
    dominate, consistent with Thompson et al.\ (2013).}
    \label{fig:varimp}
\end{figure}

\paragraph{MDS Proximity Plot.}
The MDS plot (Figure~\ref{fig:mds}) visualizes the proximity structure of the SMOTE +
Tomek training data based on Random Forest similarities. MDS1 explains 42.1\% of the
variance, clearly separating League~1 (red) on the left from higher leagues on the right.
Leagues~5 and~6 (blue and purple) form a tight cluster on the right, indicating they are
most similar to each other. Mid-tier leagues (2--4) overlap substantially in the center,
consistent with the difficulty in distinguishing them observed in the confusion matrices.
The SMOTE + Tomek model is selected for this visualization because its higher minority
recall (0.6667) produces a clearer proximity structure, better illustrating the
separation between League~1 and higher leagues.

\begin{figure}[H]
    \centering
    \includegraphics[width=0.7\textwidth]{figures/mds_plot.png}
    \caption{MDS proximity plot from Random Forest. League~1 separates clearly, while
    mid-tier leagues (2--4) show substantial overlap.}
    \label{fig:mds}
\end{figure}

\subsection{OOB Error Convergence}

The OOB error stabilizes by approximately 300--400 trees (Figure~\ref{fig:oob}), with most
per-class errors converging by that point. League~6 converges earliest ($\sim$100 trees),
while League~3 shows the most variability. The choice of $\texttt{ntree}=1000$ ensures all
classes are well-converged.

\begin{figure}[H]
    \centering
    \includegraphics[width=0.7\textwidth]{figures/oob_error.png}
    \caption{OOB error convergence by class. All classes stabilize by 400 trees.}
    \label{fig:oob}
\end{figure}

\subsection{Pairwise Binary Classification}

To enable direct comparison with Thompson et al.\ (2013), we replicated their pairwise
binary classification setup using RF + Class Weights (Table~\ref{tab:pairwise}). Our model
outperforms their Conditional Inference Forest in 3 out of 4 pairs. The only exception is
Gold--Diamond, where our model scores 75.00\% compared to their 79.01\%, likely
reflecting the inherent difficulty of distinguishing mid-tier leagues, where behavioral
differences between players are less pronounced.

Detailed discussion of these results in relation to Thompson et al.\ is provided in
Section~8.3.

\begin{table}[H]
\centering
\caption{Pairwise binary classification: our RF + Class Weights vs.\ Thompson et al.'s
Conditional Inference Forest.}
\label{tab:pairwise}
\begin{tabular}{lccc}
\toprule
\textbf{League Pair} & \textbf{Thompson et al.} & \textbf{Ours (RF + CW)} & \textbf{Outperforms?} \\
\midrule
Bronze--Gold         & 82.32\% & \textbf{88.19\%} & Yes \\
Silver--Platinum     & 78.27\% & \textbf{81.90\%} & Yes \\
Gold--Diamond        & \textbf{79.01\%} & 75.00\% & No  \\
Platinum--Masters    & 80.70\% & \textbf{84.67\%} & Yes \\
\bottomrule
\end{tabular}
\end{table}

\textit{Note:} Both classifiers are trained specifically on each league pair.
The primary difference lies in the algorithm: Conditional Inference Forest
(Thompson et al.) versus Random Forest with Class Weights, our best-performing
configuration (this study).


% ══════════════════════════════════════════════════════════
\section{Discussion}
% ══════════════════════════════════════════════════════════

\subsection{Imbalance Strategy Comparison}

\textbf{Baseline} preserves the original data distribution and achieves the highest raw
accuracy, but largely ignores the minority class. \textbf{Class Weights} provides a
meaningful minority recall improvement with only marginal accuracy cost, without modifying
the training data. However, it does not address boundary overlap between adjacent leagues.
\textbf{SMOTE + Tomek} produces the highest minority recall across all models, but synthetic
samples may not reflect real player behavior, and Tomek removal discards useful training
information, leading to the largest accuracy drops. Among the three, Class Weights offers
the best tradeoff and is selected as the preferred strategy.

\subsection{Multinomial vs.\ Pairwise Classification}

Pairwise binary classification simplifies the problem by isolating two leagues at a time,
naturally yielding higher accuracy. Multinomial classification (predicting all six leagues
simultaneously) is fundamentally harder, as adjacent leagues share similar behavioral
profiles and the model must learn all boundaries at once. In the context of a 6-class
problem where random chance yields only 16.7\%, the accuracy range of 0.28--0.41 is
reasonable. Class Weights provides the most meaningful improvement, balancing minority
recall gains without sacrificing overall discriminative power.

\subsection{Comparison with Thompson et al.}

Thompson et al.\ (2013) applied Conditional Inference Forest to pairwise binary
classification across adjacent leagues, achieving accuracies ranging from 78.27\% to
82.32\%. Using RF + Class Weights, we replicated this setup and outperformed their
results in 3 out of 4 pairs.

The only exception is Gold--Diamond (League~3 vs~5), where our model scores 75.00\%
compared to their 79.01\%. Mid-tier leagues share the most similar behavioral profiles,
making this boundary the hardest to separate regardless of model configuration.

It should also be noted that direct comparison has limitations. Their study framed
the problem as a series of binary decisions, while this study addresses 6-class
multinomial classification, a fundamentally harder task. The pairwise results here
serve as a reference point rather than a strict benchmark.

\subsection{Limitations \& Future Work}

Leagues~7 and~8 were excluded due to data quality issues, limiting scope. The relatively
small training set constrains model performance, particularly for minority classes.

Future directions include: ordinal logistic regression to capture the inherent ordering of
league ranks; XGBoost or gradient boosting methods for stronger performance on imbalanced
tabular data; and longitudinal analysis tracking the same players over time to reveal how
behavioral patterns evolve with skill development.

% ══════════════════════════════════════════════════════════
% References
% ══════════════════════════════════════════════════════════
\bibliographystyle{apalike}

\begin{thebibliography}{9}

\bibitem[Thompson et~al., 2013]{thompson2013}
Thompson, J.~J., Blair, M.~R., Chen, L., \& Henrey, A.~J. (2013).
Video game telemetry as a critical tool in the study of complex skill learning.
\textit{PLOS ONE, 8}(9), e75129.
\url{https://doi.org/10.1371/journal.pone.0075129}

\bibitem[Blair et~al., 2013]{blair2013skillcraft}
Blair, M., Thompson, J., Henrey, A., \& Chen, B. (2013).
SkillCraft1 Master Table Dataset [Dataset].
UCI Machine Learning Repository.
\url{https://doi.org/10.24432/C5161N}

\end{thebibliography}


% ══════════════════════════════════════════════════════════
\appendix
\section{Code Availability}
% ══════════════════════════════════════════════════════════

The complete R code for this analysis is available at:\\
\url{https://github.com/ShengPeiWilliam/skillcraft-ml}

\end{document}